\label{sec:patching}

This section addresses the two non-negotiable referee concerns that arise whenever an ``analytic patch'' is used inside a constructive RG:
\begin{enumerate}[label=(\roman*)]
\item the \emph{physical} lattice state must remain the unmodified Wilson Yang--Mills measure, hence reflection positivity is available \emph{a priori};
\item any patched/truncated RG step must be shown to be an \emph{analytic surrogate} that does not change the continuum limit of local/flowed observables.
\end{enumerate}

Throughout this paper we implement the following discipline.

\begin{itemize}
\item The lattice states $\omega_a$ are always the \emph{unpatched} Wilson Yang--Mills expectations (with gauge fixing used only for coordinates, never for the measure).
\item ``Patched'' refers to a \emph{chart-truncated one-step integration operator} used inside the dyadic RG to obtain uniform analytic constants.
The patched operator is never used to \emph{define} the physical lattice state.
\item De-patching is proved at the level actually used later: after exact Wilson recursion is identified, the finite-depth difference between unpatched Wilson evaluations and patched/chart-truncated RG evaluations of local or localized flowed observables is $O(u^N)$ for all $N$ on the AF ball; consequently the continuum limits computed with the patched RG agree with the physical Wilson limits.
\end{itemize}

Let $\Lambda\subset a\mathbb Z^4$ be a finite periodic lattice with a time-reflection hyperplane $\Pi=\{x_0=0\}$ and reflection map
$\theta(x_0,\mathbf x)=(-x_0,\mathbf x)$.
Let $\mathcal A_+$ denote the algebra generated by bounded cylinder functionals of link variables supported in $\{x_0>0\}$.
Let $\mathcal A^{\mathrm{gi}}_+\subset\mathcal A_+$ denote the gauge-invariant subalgebra (the physical observable sector).

\begin{remark}[Gauge averaging and enlargement of the test algebra]
Because the Wilson measure is gauge invariant, every bounded cylinder functional $F$ has the same expectation as its local gauge-average
$F^{\mathrm{gi}}:=\int F(U^g)\,dg$ (Haar product over gauge transformations on the finitely many sites that meet the support of $F$).
Thus any pointwise convergence statement proved on $\mathcal A^{\mathrm{gi}}_+$ automatically extends to $\mathcal A_+$.
We use $\mathcal A_+$ as a technically convenient test algebra (e.g.\ for quasi-local positive-time representatives of flowed probes),
while the reconstructed sharp local theory is generated by the gauge-invariant local sector.
\end{remark}

\begin{theorem}[Reflection positivity of Wilson lattice Yang--Mills]\label{thm:RP_Wilson}
For compact gauge group $G$, the Wilson lattice Yang--Mills expectation $\omega_a(\cdot)$ is reflection positive:
for every $F\in\mathcal A_+$,
\[
\omega_a(\Theta F\cdot F)\ \ge\ 0,
\]
where $\Theta$ is the Osterwalder--Schrader reflection involution (time reflection composed with $\ast$).
\end{theorem}

\begin{proof}
See Appendix~\ref{app:RP_Wilson} (Theorem~\ref{thm:RP_Wilson_full}); for background on OS reflection positivity and its lattice gauge theory implementation, see e.g.\ \cite{OS1,OS2,OSSeiler1978}.
\end{proof}

Reflection positivity implies a family of ``chessboard'' estimates that upgrade one-block bounds to bounds for translated/reflected tilings.
We record the version needed for Lemma~\ref{lem:local_tail}.

\begin{lemma}[Chessboard estimate from reflection positivity]\label{lem:chessboard}
Let $\omega$ be a reflection-positive, translation-invariant lattice state on a periodic lattice of even side length, with reflections across coordinate hyperplanes.
Let $E$ be an event (indicator functional) depending only on variables in a single fundamental reflection block $B$.
Let $E^{\#}$ denote the event obtained by reflecting/translating $E$ to all blocks in the standard chessboard tiling, and let $N_\#$ be the number of blocks in the tiling.
Then
\begin{equation}\label{eq:chessboard_general}
\omega(E)\ \le\ \big(\omega(E^{\#})\big)^{1/N_\#}.
\end{equation}
\end{lemma}

\begin{proof}
See Appendix~\ref{app:RP_Wilson} (Theorem~\ref{thm:RP_Wilson_full}); for background on OS reflection positivity and its lattice gauge theory implementation, see e.g.\ \cite{OS1,OS2,OSSeiler1978}.
\end{proof}

\begin{lemma}[Reflected OS matrix elements are stable under tuned limits]\label{lem:RP_limit}
Let $\{\omega_n\}_{n\ge 1}$ be reflection-positive states on $(\mathcal A_+,\Theta)$, and let $\mathbb D\subset[0,\infty)$ contain $0$.
Assume that for every bounded cylinder observables $F,G\in\mathcal A_+$ and every $s\in\mathbb D$ for which $\tau_s G$ is defined on the $n$th lattice, the limit
\[
K_s(F,G):=\lim_{n\to\infty}\omega_n\bigl(\Theta F\,\tau_s G\bigr)
\]
exists.
Then $K_0$ is a positive semidefinite sesquilinear form on bounded positive-time cylinders: for every bounded cylinder $F\in\mathcal A_+$,
\[
K_0(F,F)\ge 0.
\]
If in addition $\omega(F):=\lim_{n\to\infty}\omega_n(F)$ exists for every bounded cylinder $F\in\mathcal A_+$, then $\omega$ is a state on the bounded positive-time algebra and, on bounded cylinders,
\[
\omega(\Theta F\,G):=K_0(F,G).
\]
\end{lemma}

\begin{proof}
For each bounded cylinder $F\in\mathcal A_+$, reflection positivity of $\omega_n$ gives
\[
\omega_n(\Theta F\,F)\ge 0\qquad(n\ge 1).
\]
Taking $n\to\infty$ yields $K_0(F,F)\ge 0$.
Sesquilinearity of $K_0$ is inherited from the linearity of the finite-$n$ functionals together with pointwise convergence on finite linear combinations of bounded cylinders.
If the one-point limits $\omega(F)$ also exist, then $\omega$ is a state on the bounded positive-time algebra by pointwise limits of the states $\omega_n$, and the last display simply records the notation induced by the limiting reflected pairings.
\end{proof}

Lane~C uses a chart-truncated one-step integration operator to obtain uniform convexity and cumulant bounds on a fixed exponential chart.
We record the only de-patching estimate needed for the continuum limit.

\begin{lemma}[Conditioning/truncation error for bounded observables]\label{lem:depatch_block}
Let $(\Omega,\mathscr F,\mu)$ be a probability space and let $E\in\mathscr F$ satisfy $\mu(E)>0$.
Define the conditioned measure
\[
\mu^{E}(d\omega):=\frac{\mathbf 1_{E}(\omega)}{\mu(E)}\,\mu(d\omega).
\]
Then for any bounded measurable $F$ with $\|F\|_\infty\le M$,
\[
\big|\mathbb E_{\mu}[F]-\mathbb E_{\mu^{E}}[F]\big|
\ \le\ 2M\,\mu(E^c).
\]
\end{lemma}

\begin{proof}
Write $\mathbb E_{\mu^{E}}[F]=\mathbb E_\mu[F\,\mathbf 1_{E}]/\mu(E)$.
Then
\[
\mathbb E_{\mu}[F]-\mathbb E_{\mu^{E}}[F]
=\mathbb E_\mu[F\,\mathbf 1_{E^c}]
+\Bigl(1-\frac{1}{\mu(E)}\Bigr)\mathbb E_\mu[F\,\mathbf 1_{E}].
\]
Taking absolute values and using $|1-1/\mu(E)|=\mu(E^c)/\mu(E)$ gives
\[
\big|\mathbb E_{\mu}[F]-\mathbb E_{\mu^{E}}[F]\big|
\le M\,\mu(E^c)+\frac{\mu(E^c)}{\mu(E)}\,M\,\mu(E)
=2M\,\mu(E^c).
\]
\end{proof}

\paragraph{How Lemma~\ref{lem:depatch_block} is used.}
We apply Lemma~\ref{lem:depatch_block} both to a single block fluctuation law and, when needed, to the product law on a finite union of blocks,
with $E$ equal to the corresponding conjunction of good-block events.
This is the basic measure-theoretic input behind every one-step de-patching estimate below.

\paragraph{What the global tail lemma does -- and does not do.}
Lemma~\ref{lem:local_tail} is a finite-plaquette bound in the \emph{physical Wilson state}.
It is the correct large-$\beta$ input for the later Wilson-to-chart-truncated transport argument once the exact recursion has been identified, because the bad-event contribution is then averaged back in the Wilson state.
We do \emph{not} use Lemma~\ref{lem:local_tail} to assert a boundary-uniform pointwise estimate for the exact one-step conditional block factor.

\begin{lemma}[Local plaquette tail bound from chessboard + small-ball lower bound]\label{lem:local_tail}
Fix $\varepsilon\in(0,1]$ and consider the Wilson lattice gauge measure $\mu_\beta$ at coupling $\beta>0$ on a periodic $4D$ torus with compact gauge group $G$.
Let $p$ be any plaquette and write
\[
\theta(U_p):=\mathrm{dist}_{G}(U_p,\mathbf 1)
\]
for the intrinsic distance on $G$ induced by the fixed inner product on $\mathfrak g$.
Let $d_G:=\dim(G)$.

Then there exist constants
\[
\varepsilon_{\mathrm{tail}}(G)\in(0,1],\qquad c_{\mathrm{tail}}(G)>0,\qquad C_{\mathrm{tail}}(G)\ge 1,
\]
explicit in principle and depending only on $(G,\rho)$ and the fixed lattice dimension $d=4$, such that for all $\varepsilon\in(0,\varepsilon_{\mathrm{tail}}(G)]$,
\begin{equation}\label{eq:one_plaq_tail}
\mu_\beta\big(\theta(U_p)>\varepsilon\big)\ \le\ C_{\mathrm{tail}}(G)\;\varepsilon^{-2d_G/3}\;
\exp\!\Big(-c_{\mathrm{tail}}(G)\,\beta\,\varepsilon^2\Big).
\end{equation}
More generally, for any finite plaquette set $\Lambda$,
\begin{equation}\label{eq:finite_set_tail}
\mu_\beta\big(G_{\Lambda}(\varepsilon)^c\big)\ \le\ |\Lambda|\;C_{\mathrm{tail}}(G)\;\varepsilon^{-2d_G/3}\;
\exp\!\Big(-c_{\mathrm{tail}}(G)\,\beta\,\varepsilon^2\Big),
\end{equation}
where $G_{\Lambda}(\varepsilon):=\bigcap_{p\in\Lambda}\{\theta(U_p)\le \varepsilon\}$.
In particular, taking $\Lambda=\Lambda(\mathcal Q)$ the plaquettes in a fixed unit coarse block $\mathcal Q$ yields the required bound on
$\mu_\beta(G_{\mathcal Q}^c)$ used in Lemma~\ref{lem:depatch_block}.
\end{lemma}

\begin{proof}
Let $E:=\{\theta(U_{p_0})>\varepsilon\}$ for a reference plaquette $p_0$.
Let $E^{\#}$ denote the full chessboard tiling of $E$ (i.e.\ the event that \emph{every} plaquette in the torus satisfies $\theta(U_p)>\varepsilon$).
By the chessboard inequality (Lemma~\ref{lem:chessboard}),
\begin{equation}\label{eq:chess_tail_step}
\mu_\beta(E)\ \le\ \mu_\beta(E^{\#})^{1/|P|},
\end{equation}
where $|P|$ is the number of plaquettes.

\paragraph{Wilson action normalization.}
We write the Wilson action (associated to the fixed Wilson representation $\rho$ of the main text) in the form
\[
S_\beta(U)=\beta\sum_{p\in P}A_\rho(U_p),
\qquad
A_\rho(U):=1-\frac{1}{\dim(\rho)}\,\Re\Tr_{\rho}(U).
\]
(Up to the additive constant $\beta|P|$ this is equivalent to the more common normalization
$S_\beta= -\frac{\beta}{\dim(\rho)}\sum_p \Re\Tr_{\rho}(U_p)$.)

\paragraph{Local quadratic control and a global tail lower bound.}
By Taylor expansion at $\mathbf 1$ (Appendix~\ref{app:normalization_crosscheck}, Lemma~\ref{lem:wilson_quadratic_coefficient})
there exist $\delta_0(G)\in(0,1]$ and constants $0<m_-(G)\le m_+(G)<\infty$ such that for all $U=\exp(X)$ with $\|X\|\le \delta_0(G)$,
\[
m_-(G)\,\|X\|^2\ \le\ A_\rho(\exp X)\ \le\ m_+(G)\,\|X\|^2.
\]
Shrink $\delta_0(G)$ if needed so that $m_+(G)<4m_-(G)$.

Since $A_\rho$ is continuous and $A_\rho(U)=0$ only at $U=\mathbf 1$, the compact set $\{\theta(U)\ge \delta_0(G)\}$ has strictly positive minimum
\[
a_{\mathrm{out}}(G):=\min\{A_\rho(U):\ \theta(U)\ge \delta_0(G)\}\ >\ 0.
\]
Define
\[
\varepsilon_{\mathrm{tail}}(G):=\min\Big\{\delta_0(G),\ \sqrt{\frac{a_{\mathrm{out}}(G)}{m_-(G)}}\Big\}.
\]
Then for $\varepsilon\in(0,\varepsilon_{\mathrm{tail}}(G)]$ and any $U$ with $\theta(U)\ge\varepsilon$ we have the uniform lower bound
$A_\rho(U)\ge m_-(G)\varepsilon^2$.

\paragraph{Bad partition function upper bound.}
On $E^{\#}$, every plaquette satisfies $\theta(U_p)>\varepsilon$, hence $A_\rho(U_p)\ge m_-(G)\varepsilon^2$.
Therefore
\begin{equation}\label{eq:Z_bad_upper}
Z_\beta(E^{\#})
:=\int \mathbf 1_{E^{\#}}(U)\,e^{-S_\beta(U)}\,dU
\ \le\
\exp\!\Big(-\beta\,m_-(G)\,\varepsilon^2\,|P|\Big),
\end{equation}
using that the total Haar volume of link configurations is $1$.

\paragraph{Good partition function lower bound.}
For the partition function $Z_\beta=\int e^{-S_\beta(U)}\,dU$ we lower bound by restricting all links to a small ball around $\mathbf 1$.
Fix $\delta:=\varepsilon/8$ and let $B(\delta)\subset G$ be the geodesic ball of radius $\delta$ around $\mathbf 1$.
Restrict the link variables to $\{U_e\in B(\delta)\ \forall e\}$.
By the triangle inequality for the bi-invariant distance, each plaquette product satisfies
$\theta(U_p)\le 4\delta=\varepsilon/2$.
For $\varepsilon\le \varepsilon_{\mathrm{tail}}(G)\le\delta_0(G)$ we are within the quadratic regime, hence
\[
A_\rho(U_p)\ \le\ m_+(G)\,\theta(U_p)^2\ \le\ m_+(G)\,\varepsilon^2/4.
\]
Therefore
\begin{equation}\label{eq:Z_good_lower}
Z_\beta\ \ge\ \mu_{\mathrm{Haar}}(B(\delta))^{|E|}\;
\exp\!\Big(-\beta\,\frac{m_+(G)}{4}\,\varepsilon^2\,|P|\Big),
\end{equation}
where $|E|$ is the number of edges.

\paragraph{Small-ball Haar lower bound.}
Since Haar measure has a smooth positive density in exponential coordinates near the identity,
there exist constants $c_{\mathrm{ball}}(G)>0$ and (possibly smaller) $\delta_1(G)\in(0,1]$ such that for all $\delta\in(0,\delta_1(G)]$,
\begin{equation}\label{eq:Haar_ball_lower}
\mu_{\mathrm{Haar}}(B(\delta))\ \ge\ c_{\mathrm{ball}}(G)\,\delta^{d_G}.
\end{equation}

\paragraph{Combine.}
Combining \eqref{eq:Z_bad_upper}--\eqref{eq:Z_good_lower} gives
\[
\mu_\beta(E^{\#})
=\frac{Z_\beta(E^{\#})}{Z_\beta}
\le
\mu_{\mathrm{Haar}}(B(\delta))^{-|E|}
\exp\!\Big(-\beta\big(m_-(G)-m_+(G)/4\big)\varepsilon^2\,|P|\Big).
\]
On a periodic $4D$ torus one has $|E|/|P|=4/6=2/3$, hence
\[
\mu_\beta(E^{\#})^{1/|P|}
\le
\mu_{\mathrm{Haar}}(B(\delta))^{-2/3}\,
\exp\!\Big(-\beta\big(m_-(G)-m_+(G)/4\big)\varepsilon^2\Big).
\]
Using \eqref{eq:Haar_ball_lower} with $\delta=\varepsilon/8\le \delta_1(G)$ yields
\[
\mu_\beta(E^{\#})^{1/|P|}
\le
c_{\mathrm{ball}}(G)^{-2/3}\,8^{2d_G/3}\,\varepsilon^{-2d_G/3}\,
\exp\!\Big(-\beta\big(m_-(G)-m_+(G)/4\big)\varepsilon^2\Big).
\]
Together with \eqref{eq:chess_tail_step} this gives \eqref{eq:one_plaq_tail} with
\[
c_{\mathrm{tail}}(G):=m_-(G)-m_+(G)/4,
\qquad
C_{\mathrm{tail}}(G):=c_{\mathrm{ball}}(G)^{-2/3}\,8^{2d_G/3}.
\]
The finite-set bound \eqref{eq:finite_set_tail} follows by a union bound over $p\in\Lambda$ and translation invariance.
\end{proof}

\begin{lemma}[Flow localization reduces de-patching to finite-support observables]\label{lem:flow_to_depatch_bridge}
Fix a flow time $t\in(0,t_\ast]$ and a bounded flowed observable $F\in\mathcal A_{\mathrm{flow}}$.
For $R\ge 1$, let $F^{(R)}$ denote the truncation obtained by freezing the initial field outside the union of radius-$R\sqrt t$ neighborhoods around the finitely many spacetime regions that enter the definition of $F$.
Then $F^{(R)}$ is a bounded cylinder observable depending only on finitely many unflowed links, contained in at most $C_F(1+R)^4$ unit coarse blocks for a constant $C_F<\infty$ depending only on the geometry of $F$.
Moreover, there exist constants $C_F^{\mathrm{flow}},c_F^{\mathrm{flow}}>0$ such that
\[
\|F-F^{(R)}\|_\infty \ \le\ C_F^{\mathrm{flow}}\, t^{-5}\,e^{-c_F^{\mathrm{flow}}R^2}.
\]
Consequently, for every integer $N\ge 1$ there exists a choice
\[
R=R_N(a,F)\ \asymp\ \sqrt{\log(u(a)^{-1})}
\]
along the AF window such that
\[
2\|F-F^{(R)}\|_\infty\ \le\ C_N(F)\,u(a)^N.
\]
Hence finite-support control of the Wilson-to-patched difference reduces to the same comparison for the truncated observable $F^{(R)}$.
\end{lemma}

\begin{proof}
By definition, $\mathcal A_{\mathrm{flow}}$ is generated by finite products of smeared flowed curvature composites.
Apply the exponential-tail flow quasi-locality theorem for the flowed probe family (Companion~II, Appendix~B; source Theorem~5.12 in the Lane~A proof home) to each flowed factor entering $F$ and combine the resulting bounds with the boundedness of the remaining factors.
Because only finitely many insertions and smearing regions occur in $F$, the truncation depends on finitely many underlying links and occupies at most $C_F(1+R)^4$ unit coarse blocks.
The displayed exponential bound follows, after adjusting constants, from the same exponential-tail flow quasi-locality estimate.
Finally, along the AF window, choosing $R$ proportional to $\sqrt{\log(u(a)^{-1})}$ makes the exponential tail dominate any fixed power $u(a)^N$.
\end{proof}

\begin{proposition}[Finite-depth surrogate stability]\label{prop:depatch_finite_depth}
Fix any quantity later used in the paper that is obtained by composing at most $m$ one-step RG maps, blockwise reblocking/extraction operators, and a final evaluation map that is uniformly Lipschitz on the AF ball.
Assume $m\le C_{\mathrm{depth}}\log(1/u)$ on the asymptotically free window, as happens for the entrance depth and the observation depths used later, and assume that the intermediate observables encountered in this finite-depth architecture remain supported in at most polynomially many unit coarse blocks in the depth parameter
(for flowed observables, after the localization step of Lemma~\ref{lem:flow_to_depatch_bridge}, polynomially many also in the localization radius $R$).
Then for every integer $N\ge 1$ there exist $u_0(N)>0$ and $C_N<\infty$ such that replacing every one-step map in the chosen RG architecture by its patched/chart-truncated counterpart changes the resulting quantity by at most
\[
C_N\,u^N
\]
uniformly in the cutoff $a$ whenever $u\le u_0(N)$.
\end{proposition}

\begin{proof}
After Theorem~\ref{thm:wilson_exact_recursion}, every exact one-step map appearing in the Wilson evaluation is the conditional expectation operator $T_{j+1}$, while the patched architecture replaces it by the corresponding good-event truncation $T^{\mathrm{good}}_{j+1}$ with the same deterministic post-processing.
Write the finite-depth difference as a telescoping sum over the at most $m$ substituted one-step kernels.
At the $r$-th telescoping term one encounters a bounded intermediate observable $H_r$ supported in a finite union $\Lambda_r$ of unit coarse blocks; by hypothesis, $|\Lambda_r|$ is bounded by a fixed polynomial in $r$ (and, for flowed observables after localization, also polynomially in $R$).

Because $H_r$ depends only on the blocks in $\Lambda_r$, truncating all other block factors is immaterial.
Hence Proposition~\ref{prop:averaged_bad_event_transport} gives
\[
\big|\omega_a(T_{j+1}H_r)-\omega_a(T_{j+1}^{\mathrm{good}}H_r)\big|
\le
C(H_r)\,|\Lambda_r|\,\exp\!\big(-c_{\mathrm{tail}}(G)\beta\,\varepsilon^2\big).
\]
The deterministic reblocking/extraction operators and the final evaluation map are uniformly Lipschitz on the AF ball by Theorem~\ref{thm:RG_analytic_bounds} and Theorem~\ref{thm:C3_4}, so the telescoping sum is bounded by a constant times
\[
\Big(\sum_{r=0}^{m-1}(1+r)^{\nu}\Big)\,(1+R)^{\nu_F}\,\exp\!\big(-c_{\mathrm{tail}}(G)\beta\,\varepsilon^2\big)
\]
for some fixed exponents $\nu,\nu_F$ depending only on the chosen quantity.
Since $m\le C_{\mathrm{depth}}\log(1/u)$ and $R$ may be taken $O(\sqrt{\log(1/u)})$, the prefactor is at worst polylogarithmic in $1/u$.
Along the AF window one has $\beta\asymp u^{-1}$, so the stretched exponential $\exp(-c/u)$ dominates every polylogarithm.
Shrinking $u_0(N)$ if necessary therefore yields an overall bound $C_Nu^N$.
\end{proof}

\begin{theorem}[Support propagation for the actual finite-depth Lane~C recursion]\label{thm:laneC_support_propagation}
Fix a bounded cylinder observable $H_0$ supported in a finite union $\Lambda_0$ of unit coarse blocks at the input scale, and let $H_r$ denote the intermediate observable obtained after $r$ steps of the actual chart-truncated Lane~C recursion of Definition~\ref{def:RG_good}, before the final scalar evaluation.
Then there is a canonical dyadic parent-hull operation $\operatorname{Par}$ on finite block families such that
\[
\operatorname{supp}(H_r)\subseteq \operatorname{Par}^r(\Lambda_0)
\qquad (r\ge0),
\]
and
\[
|\operatorname{Par}^r(\Lambda_0)|\le |\Lambda_0|.
\]
In particular, for bounded cylinder observables the number of blocks touched by the actual finite-depth Lane~C recursion is uniformly bounded in the depth.
If $H_0=F^{(R)}$ is the localization of a flowed observable from Lemma~\ref{lem:flow_to_depatch_bridge}, then
\[
|\operatorname{Par}^r(\operatorname{supp}F^{(R)})|\le C_F(1+R)^4
\qquad (r\ge0),
\]
so the finite-depth locality hypothesis of Proposition~\ref{prop:depatch_finite_depth} is verified concretely for the actual Lane~C recursion.
\end{theorem}

\begin{proof}
For one actual Lane~C step at scale $j$, Proposition~\ref{prop:RGgood_support_one_step} gives the explicit operator identity
\[
\operatorname{supp}(\mathcal R_j^{\mathrm{good}}H)\subseteq \operatorname{Par}(\Lambda)
\]
whenever $H$ is supported in a finite family $\Lambda$ of $j$-scale blocks.
This is the concrete operator-by-operator support map for the actual chart-truncated fluctuation operator, deterministic reblocking/rescaling, and exact extraction/reparametrization.

Iterating this one-step inclusion starting from $H_0$ supported in $\Lambda_0$ gives
\[
\operatorname{supp}(H_r)\subseteq \operatorname{Par}^r(\Lambda_0)
\qquad (r\ge0).
\]
Because each half-open child block has a unique dyadic parent block, distinct child blocks may merge but no parent block can split; hence
\[
|\operatorname{Par}^r(\Lambda_0)|\le |\Lambda_0|.
\]
If $H_0=F^{(R)}$ is the localization of a flowed observable from Lemma~\ref{lem:flow_to_depatch_bridge}, that lemma gives
\[
|\operatorname{supp}F^{(R)}|\le C_F(1+R)^4
\]
at the input scale, and the same parent-hull monotonicity preserves the bound at every later depth:
\[
|\operatorname{Par}^r(\operatorname{supp}F^{(R)})|\le C_F(1+R)^4.
\]
This is exactly the support-growth statement required in Proposition~\ref{prop:depatch_finite_depth} for the actual Lane~C recursion.
\end{proof}

\paragraph{Architectural clarification.}
The finite-range / product structure invoked later in Lane~C is not an ad hoc surrogate ingredient.
For the Wilson theory itself, once the block-skeleton links are held fixed, the complementary interior-link law factorizes exactly over $(j+1)$-blocks.
Thus the only auxiliary modification used in the analytic RG is the \emph{local good-block chart truncation}; the earlier Wilson-to-product issue is resolved by the exact one-step factorization theorem below.

\begin{theorem}[Exact one-step Wilson block factorization]\label{thm:wilson_one_step_factorization}
Fix a dyadic scale $j$ and partition the $j$-lattice into half-open $(j+1)$-blocks of side $2$.
Let $\Sigma_{j+1}$ denote the block skeleton, i.e.\ the set of $j$-links not strictly contained in the interior of a unique half-open $(j+1)$-block, and write $U_{\Sigma}$ for the link configuration on $\Sigma_{j+1}$.
Let $U_{\mathrm{int}}$ denote the complementary interior-link variables, and let $\mathcal B_{j+1}$ denote the family of half-open $(j+1)$-blocks and, for each block $\mathcal Q\in \mathcal B_{j+1}$, let $U_{\mathcal Q}$ be the collection of interior links belonging to $\mathcal Q$.
Then the Wilson lattice gauge measure has the conditional product form
\begin{equation}\label{eq:wilson_block_factorization}
\mu^{\mathrm W}_{j}(dU_{\mathrm{int}}\mid U_{\Sigma})
\;=\;
\bigotimes_{\mathcal Q\in \mathcal B_{j+1}} \mu^{\mathrm W}_{j+1,\mathcal Q}\bigl(dU_{\mathcal Q}\mid U_{\Sigma}|_{\partial \mathcal Q}\bigr),
\end{equation}
where each block factor depends only on the skeleton data adjacent to $\mathcal Q$.
After imposing the canonical tree gauge inside $\mathcal Q$ with boundary data frozen, each block factor is represented by the finite-dimensional local density used later in Definitions~\ref{def:local_vars} and \ref{def:fluc_density}.
In particular, the exact Wilson one-step fluctuation law satisfies the finite-range input \textup{(FC1)} with $R_{\mathrm{fr}}=0$ conditional on the skeleton variables, and the chart-truncated law obtained by imposing the good-block event in each block preserves the same block-product structure.
\end{theorem}

\begin{proof}
Write the fine plaquette with anchor $x$ and directions $(\mu,\nu)$ as $p(x;\mu,\nu)$, and let $\mathcal Q(x)$ denote the unique half-open $(j+1)$-block containing the anchor $x$.
Because the blocks are half-open, every non-skeleton link of $p(x;\mu,\nu)$ is strictly contained in the same owner block $\mathcal Q(x)$; any remaining links of the plaquette lie on the block skeleton.
Hence no plaquette term can involve interior links from two distinct blocks.

The Wilson measure is the product Haar measure on links multiplied by the local Boltzmann factor
\[
\exp\Bigl\{-\beta\sum_{p}A_\rho(U_p)\Bigr\},
\qquad
A_\rho(U):=1-\frac{1}{\dim(\rho)}\,\Re\Tr_{\rho}(U).
\]
After $U_{\Sigma}$ is held fixed, the conditional action is therefore a sum of block-local terms,
\[
S_j(U_{\mathrm{int}};U_{\Sigma})=\sum_{\mathcal Q\in\mathcal B_{j+1}} S_{\mathcal Q}\bigl(U_{\mathcal Q};U_{\Sigma}|_{\partial \mathcal Q}\bigr),
\]
and the product Haar prior on $U_{\mathrm{int}}$ already factorizes over blocks.
This gives \eqref{eq:wilson_block_factorization}.
Inside each block one may then impose the canonical tree gauge with the boundary links frozen; by Haar bi-invariance on the finite block, the resulting change of variables is blockwise and produces exactly the finite-dimensional local coordinate description used in Lane~C.
Finally, the good-block event is measurable in the corresponding block factor, so truncating by the blockwise indicators preserves the product decomposition.
\end{proof}

\begin{theorem}[Nested Wilson conditioning gives the exact dyadic recursion]\label{thm:wilson_exact_recursion}
Fix a starting dyadic scale $j_0$ and a finite depth $m\ge 1$.
For each $r\in\{0,\dots,m\}$, let $\Sigma_{j_0+r}$ be the dyadic skeleton obtained after $r$ repetitions of the half-open $2$-blocking rule, and let $\mathcal F_{j_0+r}$ be the sigma-algebra generated by the links on $\Sigma_{j_0+r}$.
Then the skeleton sigma-algebras are nested,
\[
\mathcal F_{j_0+m}\subset\cdots\subset \mathcal F_{j_0+1}\subset \mathcal F_{j_0},
\]
and for each step $r=0,\dots,m-1$ the exact one-step conditional expectation onto $\mathcal F_{j_0+r+1}$ is given by the block-product kernel of Theorem~\ref{thm:wilson_one_step_factorization} applied on the current $(j_0+r)$-scale lattice.
Consequently, for every bounded observable $H$ measurable with respect to the scale-$j_0$ field,
\begin{equation}\label{eq:iterated_wilson_tower}
\omega_a(H)
=\omega_a\Big( T_{j_0+m}\cdots T_{j_0+1}H\Big),
\end{equation}
where
\[
(T_{j+1}G)(U_{\Sigma_{j+1}})
:=\int G(U_{\Sigma_{j+1}},U_{\mathrm{int}})\,\mu^{\mathrm W}_{j}(dU_{\mathrm{int}}\mid U_{\Sigma_{j+1}})
\]
is the exact Wilson one-step conditional expectation operator.
After the deterministic reblocking/rescaling maps and the exact coordinate re-expression from \eqref{eq:RG_density_form}, this is precisely the untruncated Lane~C recursion generated by the maps $\mathcal R_j$.
Hence the chart-truncated recursion differs from the exact Wilson recursion only by replacing each $T_{j+1}$ with its blockwise good-event truncation $T^{\mathrm{good}}_{j+1}$.
\end{theorem}

\begin{proof}
Apply Theorem~\ref{thm:wilson_one_step_factorization} at the first step $j_0\to j_0+1$.
Because each $(j+2)$-block is a union of half-open $(j+1)$-blocks, the corresponding dyadic skeletons are nested, so the sigma-algebras $\mathcal F_{j_0+r}$ form a decreasing filtration.
For each $r$, Theorem~\ref{thm:wilson_one_step_factorization} on the current $(j_0+r)$-scale lattice identifies $T_{j_0+r+1}$ with the conditional expectation $\mathbb E_{\omega_a}[\,\cdot\mid \mathcal F_{j_0+r+1}]$.
Iterating conditional expectation therefore gives the tower identity \eqref{eq:iterated_wilson_tower}.

The reblocking/rescaling maps used in Lane~C are deterministic measurable maps of the surviving skeleton variables, so they do not change the underlying conditional-expectation identity.
Likewise, the passage from the exact blocked density $\mathcal Z_{j+1}$ to the coordinates $(u',K')$ is the exact re-expression described after \eqref{eq:RG_density_form}; it introduces no approximation.
Therefore the untruncated one-step Lane~C map $\mathcal R_j$ is exactly the blocked Wilson recursion written in $(u,K)$ coordinates.
By Definition~\ref{def:RG_good}, the patched recursion keeps the same deterministic post-processing and changes only the one-step conditional expectation $T_{j+1}\mapsto T^{\mathrm{good}}_{j+1}$.
\end{proof}

\begin{corollary}[Terminal blockwise conditional product and disjoint-block independence]\label{cor:wilson_terminal_block_product}
In the setting of Theorem~\ref{thm:wilson_exact_recursion}, fix the terminal depth $m$ and the corresponding family $\mathcal B_{j_0+m}$ of half-open terminal blocks.
Conditional on the terminal skeleton sigma-algebra $\mathcal F_{j_0+m}$, the fine Wilson fluctuation law factorizes over the terminal blocks:
\begin{equation}\label{eq:wilson_terminal_block_product}
\mu^{\mathrm W}_{j_0}(dU_{\mathrm{int}}\mid \mathcal F_{j_0+m})
\;=\;
\bigotimes_{\mathcal Q\in\mathcal B_{j_0+m}} \mu^{\mathrm W}_{\mathcal Q}\bigl(dU_{\mathcal Q}\mid U_{\Sigma_{j_0+m}}|_{\partial \mathcal Q}\bigr),
\end{equation}
where $U_{\mathcal Q}$ denotes the fine links strictly interior to the terminal block $\mathcal Q$.
Consequently, if $H_1,\dots,H_r$ are bounded observables such that each $H_i$ is measurable with respect to the fine links contained in a family $\mathscr S_i\subset \mathcal B_{j_0+m}$ and the families $\mathscr S_1,\dots,\mathscr S_r$ are pairwise disjoint, then
\begin{equation}\label{eq:wilson_terminal_conditional_independence}
\mathbb E_{\omega_a}\!\Big[\prod_{i=1}^r H_i\,\Big|\,\mathcal F_{j_0+m}\Big]
\;=\;
\prod_{i=1}^r \mathbb E_{\omega_a}[H_i\mid \mathcal F_{j_0+m}].
\end{equation}
\end{corollary}

\begin{proof}
Theorem~\ref{thm:wilson_one_step_factorization} gives exact conditional block-factorization at one dyadic step.
Iterating that factorization through the nested filtration of Theorem~\ref{thm:wilson_exact_recursion} yields an exact product decomposition of the fine interior-link law over the terminal half-open blocks once the terminal skeleton field is frozen, which is \eqref{eq:wilson_terminal_block_product}.
If each $H_i$ depends only on fine links from a block family $\mathscr S_i$ and the families are pairwise disjoint, then under the product law \eqref{eq:wilson_terminal_block_product} the random variables $H_1,\dots,H_r$ are conditionally independent given $\mathcal F_{j_0+m}$.
This is exactly \eqref{eq:wilson_terminal_conditional_independence}.
\end{proof}

\begin{proposition}[Averaged one-step bad-event transport]\label{prop:averaged_bad_event_transport}
Fix a dyadic step $j\to j+1$ and let $H$ be a bounded observable depending only on the fluctuation variables in a finite union $\Lambda$ of unit coarse $(j+1)$-blocks.
Let
\[
G_{\Lambda}(\varepsilon):=\bigcap_{\mathcal Q\subset\Lambda} G_{\mathcal Q}(\varepsilon)
\]
denote the conjunction of the good-block events on those blocks, where $G_{\mathcal Q}(\varepsilon)$ is the good-block event of Definition~\ref{def:good_block}, and let $T_{j+1}^{\mathrm{good},\Lambda}$ be the one-step operator obtained from $T_{j+1}$ by conditioning only the block factors for $\mathcal Q\subset\Lambda$ on $G_{\mathcal Q}(\varepsilon)$ (leaving all other blocks exact).
Then
\begin{equation}\label{eq:avg_one_step_transport}
\big|\omega_a(T_{j+1}H)-\omega_a(T_{j+1}^{\mathrm{good},\Lambda}H)\big|
\le
2\|H\|_{\infty}\,\omega_a\!\big(G_{\Lambda}(\varepsilon)^c\big).
\end{equation}
Consequently, if $P(\Lambda)$ denotes the finite plaquette set contained in the blocks of $\Lambda$ and $d_G:=\dim(G)$, then
\begin{equation}\label{eq:avg_one_step_transport_explicit}
\big|\omega_a(T_{j+1}H)-\omega_a(T_{j+1}^{\mathrm{good},\Lambda}H)\big|
\le
2\|H\|_{\infty}\,|P(\Lambda)|\,C_{\mathrm{tail}}(G)\,\varepsilon^{-2d_G/3}\exp\!\big(-c_{\mathrm{tail}}(G)\beta\,\varepsilon^2\big).
\end{equation}
In particular, on the AF window, for every integer $N\ge1$ there exists $u_0(N,\Lambda,H)>0$ such that
\[
u\le u_0(N,\Lambda,H)\quad\Longrightarrow\quad
\big|\omega_a(T_{j+1}H)-\omega_a(T_{j+1}^{\mathrm{good},\Lambda}H)\big|
\le C_N(\Lambda,H)\,u^N.
\]
\end{proposition}

\begin{proof}
Condition on the skeleton sigma-algebra $\mathcal F_{j+1}$.
By Theorem~\ref{thm:wilson_one_step_factorization}, the exact conditional fluctuation law factorizes over the $(j+1)$-blocks, and on the blocks of $\Lambda$ the patched law is obtained by conditioning that product law on $G_{\Lambda}(\varepsilon)$.
For each fixed skeleton configuration, apply Lemma~\ref{lem:depatch_block} to the product fluctuation law on the blocks of $\Lambda$ with event $G_{\Lambda}(\varepsilon)$ to obtain
\[
\big|T_{j+1}H-T_{j+1}^{\mathrm{good},\Lambda}H\big|
\le
2\|H\|_{\infty}\,
\mu^{\mathrm W}_{j}\!\big(G_{\Lambda}(\varepsilon)^c\mid \mathcal F_{j+1}\big).
\]
Now average in the Wilson state and use the tower property:
\[
\omega_a\!\Big[\mu^{\mathrm W}_{j}\!\big(G_{\Lambda}(\varepsilon)^c\mid \mathcal F_{j+1}\big)\Big]
=
\omega_a\!\big(G_{\Lambda}(\varepsilon)^c\big).
\]
This proves \eqref{eq:avg_one_step_transport}.
The explicit estimate \eqref{eq:avg_one_step_transport_explicit} follows by applying Lemma~\ref{lem:local_tail} to the finite plaquette set $P(\Lambda)$ and using a union bound.
Since along the AF window one has $\beta\asymp u^{-1}$, the exponential term is superpolynomial in $u^{-1}$ and therefore dominates every fixed power $u^N$.
\end{proof}

\begin{proposition}[Local transport closure for the actual chart-truncated Lane~C recursion]\label{prop:local_transport_closure}
Fix a bounded cylinder observable $H$ supported in a finite union $\Lambda$ of unit coarse blocks.
Let $E_a^{\mathrm W}(H)$ denote its evaluation by the exact finite-depth Wilson recursion of Theorem~\ref{thm:wilson_exact_recursion}, and let $E_a^{\mathrm{good}}(H)$ denote its evaluation by the actual chart-truncated Lane~C recursion, i.e. by the dyadic RG architecture of Definition~\ref{def:RG_good} in which each one-step exact conditional expectation $T_{j+1}$ is replaced by the chart-truncated operator $T_{j+1}^{\mathrm{good}}$ while all deterministic reblocking/extraction maps are left unchanged.
Then for every integer $N\ge 1$ there exist $u_0(N)>0$ and $C_N(H)<\infty$ such that
\[
\big|E_a^{\mathrm W}(H)-E_a^{\mathrm{good}}(H)\big|\le C_N(H)\,u(a)^N
\]
uniformly for $a$ in the AF window.
Moreover, if $F\in\mathcal A_{\mathrm{flow}}$ is fixed and $F^{(R)}$ is the cylinder truncation from Lemma~\ref{lem:flow_to_depatch_bridge}, then for every $N$ one may choose $R=R_N(a,F)\asymp \sqrt{\log(u(a)^{-1})}$ so that
\[
\big|\omega_a(F)-\omega_a^{\mathrm{RG}}(F)\big|
\le
\big|E_a^{\mathrm W}(F^{(R)})-E_a^{\mathrm{good}}(F^{(R)})\big| + C_N(F)\,u(a)^N,
\]
and the first term is again $O(u(a)^N)$.
\end{proposition}

\begin{proof}
For bounded cylinder observables, Theorem~\ref{thm:wilson_exact_recursion} identifies the exact Wilson evaluation with a finite-depth composition of the operators $T_{j+1}$ followed by the deterministic reblocking/extraction maps of Lane~C.
By Definition~\ref{def:RG_good} and the construction in Section~\ref{sec:rg}, the actual analytic Lane~C recursion uses those same deterministic maps and changes only the one-step fluctuation integration operators by the substitutions $T_{j+1}\mapsto T_{j+1}^{\mathrm{good}}$.
Thus the one-step difference between the physical Wilson evaluation and the actual chart-truncated Lane~C evaluation is exactly the local good-block truncation treated in Proposition~\ref{prop:averaged_bad_event_transport}.
Theorem~\ref{thm:laneC_support_propagation} verifies that the actual chart-truncated Lane~C recursion keeps the intermediate supports inside a uniformly bounded parent hull (and, for flowed observables after localization, inside at most $C_F(1+R)^4$ blocks).
Therefore Proposition~\ref{prop:depatch_finite_depth} applies to the concrete Lane~C recursion and propagates the one-step averaged estimate through the full finite-depth evolution, yielding the stated $O(u^N)$ bound.
For flowed observables, apply Lemma~\ref{lem:flow_to_depatch_bridge} and choose the localization radius so that the deterministic tail from that lemma is itself $O(u(a)^N)$; then apply the cylinder estimate to $F^{(R)}$.
\end{proof}

\begin{remark}[Audit reference for the Wilson-versus-patched seam]\label{rem:transport_status}
Theorem~\ref{thm:Wilson_patch_transport} below is the single theorem-level location where the patched Lane~C estimates are returned to the physical Wilson theory.
Theorem~\ref{thm:wilson_one_step_factorization}, Theorem~\ref{thm:wilson_exact_recursion}, Proposition~\ref{prop:averaged_bad_event_transport}, and Proposition~\ref{prop:laneC_transport_identification} serve only as proof inputs to that fixed-scale closure statement.
\end{remark}

\begin{remark}[Phase-V Route~1 measure pipeline ledger]\label{rem:route1_measure_pipeline}
The exact Wilson/Lane~C measure-side chain used later is only the following theorem-level packet:
\begin{enumerate}[label=\textup{(\arabic*)},leftmargin=2.35em]
\item Theorem~\ref{thm:wilson_one_step_factorization} identifies the exact Wilson one-step fluctuation law simultaneously as conditional expectation onto the next skeleton sigma-algebra and as an exact block-product law over the half-open coarse blocks.
\item Theorem~\ref{thm:wilson_exact_recursion} iterates those exact one-step conditional expectations through the dyadic filtration and shows that the untruncated Lane~C recursion is exactly the blocked Wilson recursion written in $(u,K)$ coordinates.
\item Proposition~\ref{prop:averaged_bad_event_transport} quantifies the only local measure change later used: conditioning the relevant block factors on the good-block event over one fixed support hull.
\item Proposition~\ref{prop:laneC_transport_identification} states that the actual chart-truncated Lane~C recursion makes precisely that truncation-only substitution and no other measure-side modification.
\item Theorem~\ref{thm:Wilson_patch_transport} packages items~\textup{(1)}--\textup{(4)} into the same-scale Wilson-to-patched return used later by the public Route~1 spine.
\end{enumerate}
No later Route~1 theorem reopens the measure-side comparison below Theorem~\ref{thm:Wilson_patch_transport}; downstream uses consume only that packaged theorem-facing return.
\end{remark}

\begin{proposition}[The actual Lane~C chart truncation matches the transport template]\label{prop:laneC_transport_identification}
Fix the dyadic RG architecture of Lane~C, with chart-truncated one-step map $\mathcal R_j^{\mathrm{good}}$ from Definition~\ref{def:RG_good}.
Then, relative to the exact Wilson recursion identified in Theorem~\ref{thm:wilson_exact_recursion}, the only modification made by the Lane~C analytic recursion is the blockwise good-event truncation
\[
T_{j+1}\mapsto T_{j+1}^{\mathrm{good}}
\]
at each one-step fluctuation integration, followed by the same deterministic reblocking/rescaling and exact extraction maps used in the untruncated recursion.
Consequently the actual Lane~C chart-truncated recursion satisfies the hypotheses of Proposition~\ref{prop:local_transport_closure}.
\end{proposition}

\begin{proof}
Theorem~\ref{thm:wilson_exact_recursion} identifies the untruncated Lane~C recursion with the exact finite-depth blocked Wilson recursion written in $(u,K)$ coordinates.
Definition~\ref{def:RG_good} then shows that the chart-truncated map keeps all deterministic post-processing unchanged and alters only the one-step fluctuation expectation by conditioning each block factor on the corresponding good-block event.
Because the good-block event factorizes blockwise (Definition~\ref{def:trunc_measure}) and the exact one-step Wilson law already has block-product structure by Theorem~\ref{thm:wilson_one_step_factorization}, this is precisely the truncation-only setting abstracted in Proposition~\ref{prop:local_transport_closure}.
\end{proof}

\begin{theorem}[Same-scale Wilson-to-patched transport]\label{thm:Wilson_patch_transport}
Fix a support budget $M_{\mathrm{supp}}\in\mathbb N$ and a depth-growth constant $C_{\mathrm{depth}}>0$.
Use the fixed-scale test norm
\[
\|H\|_{\mathrm{test}}:=\|H\|_{L^\infty}
\]
on bounded cylinder observables whose input-scale support is contained in a union of at most $M_{\mathrm{supp}}$ unit coarse blocks.
For such an observable $H$, let $E_{a,m}^{\mathrm W}(H)$ denote the exact Wilson evaluation obtained by the $m$-step recursion of Theorem~\ref{thm:wilson_exact_recursion}, and let $E_{a,m}^{\mathrm{good}}(H)$ denote the parallel $m$-step chart-truncated evaluation obtained from Definition~\ref{def:RG_good}.
On the finite-volume/admissible-boundary-condition class used later in Route~1, the same notation denotes the identical local construction with the outer boundary data frozen.
Then for every integer $N\ge 1$ there exist $u_0(N,C_{\mathrm{depth}},M_{\mathrm{supp}})>0$ and $C_N(C_{\mathrm{depth}},M_{\mathrm{supp}})<\infty$, depending only on the depth-growth constant, the support budget, and the group ledger, such that for all $u(a)\le u_0$ and all dyadic depths satisfying
\[
1\le m\le C_{\mathrm{depth}}\log\!\bigl(1/u(a)\bigr),
\]
the following hold uniformly in the cutoff $a$, uniformly in the admissible depths $m$, and uniformly over the Route~1 finite-volume/admissible-boundary-condition class.
\begin{enumerate}[label=\textup{(\alph*)}]
\item For every bounded cylinder observable $H$ in the fixed-scale test class,
\[
\big|E_{a,m}^{\mathrm W}(H)-E_{a,m}^{\mathrm{good}}(H)\big|
\le C_N\,\|H\|_{\mathrm{test}}\,u(a)^N.
\]
\item For every bounded cylinder pair $H_1,H_2$ whose common input-scale support hull is contained in a union of at most $M_{\mathrm{supp}}$ unit coarse blocks (equivalently, for which $H_1$, $H_2$, and $H_1H_2$ all belong to the fixed-scale test class),
\[
\big|\operatorname{Cov}_{a,m}^{\mathrm W}(H_1,H_2)-\operatorname{Cov}_{a,m}^{\mathrm{good}}(H_1,H_2)\big|
\le C_N'\,\|H_1\|_{\mathrm{test}}\,\|H_2\|_{\mathrm{test}}\,u(a)^N,
\]
where
\[
\operatorname{Cov}_{a,m}^{\star}(H_1,H_2):=E_{a,m}^{\star}(H_1H_2)-E_{a,m}^{\star}(H_1)E_{a,m}^{\star}(H_2),\qquad \star\in\{\mathrm W,\mathrm{good}\}.
\]
In particular, if $F,G\in\mathcal A_+$ are bounded positive-time cylinder observables and the reflected/translated product $(\Theta F)\,\tau_s G$ remains in the same fixed-scale test class, then
\[
\big|E_{a,m}^{\mathrm W}((\Theta F)\tau_s G)-E_{a,m}^{\mathrm{good}}((\Theta F)\tau_s G)\big|
\le C_N''\,\|F\|_{\infty}\,\|G\|_{\infty}\,u(a)^N.
\]
If $H_1,H_2$ are centered with respect to either evaluation, this is the corresponding centered-correlation estimate.
\item For every fixed flowed local observable $F\in\mathcal A_{\mathrm{flow}}$ and every integer $N\ge 1$ there exists $C_N(F)<\infty$ such that
\[
\big|\omega_a(F)-\omega_a^{\mathrm{RG}}(F)\big|\le C_N(F)\,u(a)^N
\]
uniformly for $a$ in the AF window, where $\omega_a^{\mathrm{RG}}$ denotes the finite-$a$ functional produced by the chart-truncated dyadic RG recursion.
\end{enumerate}
\end{theorem}

\begin{proof}
Part~\textup{(a)} packages the exact Wilson recursion, the actual chart-truncated recursion, and the averaged bad-event transport estimate into one finite-depth argument with logarithmic-depth uniformity.
Theorem~\ref{thm:wilson_exact_recursion} identifies $E_{a,m}^{\mathrm W}$ as the depth-$m$ composition of the exact one-step conditional expectations $T_{j+1}$.
Proposition~\ref{prop:laneC_transport_identification} identifies $E_{a,m}^{\mathrm{good}}$ as the parallel depth-$m$ recursion in which the same deterministic reblocking/extraction maps are kept and only the substitutions $T_{j+1}\mapsto T_{j+1}^{\mathrm{good}}$ are made.
For an input observable $H$ supported in at most $M_{\mathrm{supp}}$ unit coarse blocks, Theorem~\ref{thm:laneC_support_propagation} shows that every intermediate observable encountered along the patched recursion stays supported in at most $M_{\mathrm{supp}}$ blocks.
Hence the actual Lane~C recursion falls under Proposition~\ref{prop:depatch_finite_depth} with support budget $M_{\mathrm{supp}}$ and with depth bound $m\le C_{\mathrm{depth}}\log(1/u(a))$.
That proposition already absorbs the telescoping sum, the logarithmic depth, and the polylogarithmic prefactors, and yields a bound
\[
\big|E_{a,m}^{\mathrm W}(H)-E_{a,m}^{\mathrm{good}}(H)\big|
\le C_N(C_{\mathrm{depth}},M_{\mathrm{supp}})\,\|H\|_{\infty}\,u(a)^N
\]
for every prescribed $N$ once $u(a)$ is sufficiently small.
Because the argument is local on the support hull and the deterministic support map is unchanged, the same bound is uniform over the finite-volume/admissible-boundary-condition class used in Route~1: boundary data enter only as frozen exterior variables on that hull.
This proves part~\textup{(a)}.

For part~\textup{(b)}, apply part~\textup{(a)} to $H_1$, $H_2$, and $H_1H_2$; this is legitimate because the common support-hull hypothesis in part~\textup{(b)} puts all three observables in the same fixed-scale test class. Use then the identity
\[
AB-A'B'=(A-A')B+A'(B-B')
\]
with $A=E_{a,m}^{\mathrm W}(H_1)$, $B=E_{a,m}^{\mathrm W}(H_2)$, $A'=E_{a,m}^{\mathrm{good}}(H_1)$, and $B'=E_{a,m}^{\mathrm{good}}(H_2)$.
The positive-time pairing is the special case $H=(\Theta F)\tau_s G$.
If the observables are centered for either evaluation, the covariance identity reduces immediately to the corresponding centered correlation bound.

For part~\textup{(c)}, apply Lemma~\ref{lem:flow_to_depatch_bridge} and choose the localization radius so that the deterministic truncation tail is $O(u(a)^N)$; then apply Proposition~\ref{prop:local_transport_closure} to the cylinder truncation $F^{(R)}$. The support bound from Lemma~\ref{lem:flow_to_depatch_bridge} is polynomial in $R$, so the proposition applies directly to this localized flowed observable.
\end{proof}

\begin{remark}[Phase-IV uniformity note for Theorem~\ref{thm:Wilson_patch_transport}]
The theorem-facing uniformity is exactly logarithmic in the dyadic depth: $m\le C_{\mathrm{depth}}\log(1/u(a))$. This is the regime later realized by one-shot entrance, fixed physical comparison scales, and flowed-localization closures on the public Route~1 spine. The statement should therefore neither be weakened to a merely fixed-depth estimate nor overread as uniform for arbitrary depth families outside that logarithmic window.
\end{remark}

\begin{lemma}[De-patching at the level of continuum limits]\label{lem:depatch_continuum}
Assume the chart-truncated RG outputs converge to a limiting functional $\omega_{\mathrm{RG}}$ on the flowed algebra $\mathcal A_{\mathrm{flow}}$.
Then the unpatched Wilson expectations converge with the same limit:
\[
\lim_{a\to 0}\omega_a(F)=\omega_{\mathrm{RG}}(F),\qquad F\in\mathcal A_{\mathrm{flow}}.
\]
\end{lemma}

\begin{proof}
By assumption, $\omega_a^{\mathrm{RG}}(F)\to \omega_{\mathrm{RG}}(F)$ for each fixed $F$.
Theorem~\ref{thm:Wilson_patch_transport} gives $|\omega_a(F)-\omega_a^{\mathrm{RG}}(F)|\to 0$ as $a\to0$.
Therefore $\omega_a(F)\to \omega_{\mathrm{RG}}(F)$.
\end{proof}

\paragraph{Consequence.}
On the tuned dyadic trajectory used on the load-bearing Route~1/Lane~B path, Theorem~\ref{thm:continuum_tightness} records the actual bounded positive-time Wilson limit and its reflected OS matrix elements directly.
Lemma~\ref{lem:depatch_continuum} identifies the chart-truncated RG continuum state with the Wilson continuum expectations on flowed observables in $\mathcal A_{\mathrm{flow}}$.
Thus the de-patching module is used only for the flowed/Lane~C continuum construction; reflection positivity of the bounded base state is handled separately in Theorem~\ref{thm:continuum_tightness}.
